\begin{textsample}{POS Dim 1 – sp – Score 28.00 – sp/sp\_ef\_ciencias\_humanas\_his\_1.txt}  \label{ex:f1_pos_165}
História O saber histórico na sala de aula tem se caracterizado por \textbf{um} duplo movimento.
De \textbf{um} \textbf{lado}, tenta -se compreender aspectos do presente por meio do passado¹³.
De outro, busca -se reelaborar a história a partir de novos questionamentos.
Com tal processo, pretende -se contribuir para a construção das identidades dos diferentes grupos \textbf{que} constituem a sociedade. ¹³ Desde os Annales se entende \textbf{que} o historiador é também fruto do seu tempo e \textbf{que} toda investigação histórica e, consequentemente, toda narrativa histórica é fruto de perguntas realizadas por \textbf{um} ser social \textbf{que} só as pensou graças aos paradigmas do tempo presente.
O manual \textbf{metodológico} \textbf{que} organizou essas práticas historiográficas foi o livro Apologia da História, de Marc Bloch.
Décadas depois, o historiador Michel de Certeau recupera a necessidade de considerar \textbf{que}, mesmo tratando do passado, o texto fatidicamente traz elementos do presente de quem o escreveu.
Chamado por Certeau de “ operação historiográfica ”, esse processo é a junção da “ relação entre \textbf{um} lugar ( \textbf{um} recrutamento, \textbf{um} meio, \textbf{uma} profissão, etc. ), procedimentos de análise ( \textbf{uma} \textbf{disciplina} ) e a construção de \textbf{um} texto ( \textbf{uma} literatura ) ” ( CERTEAU, 1982 ). > Os velhos marcos históricos estão sendo revistos, mesmo \textbf{que} paulatinamente, podendo -se introduzir \textbf{uma} história da Antiguidade pelas sociedades indígenas, pela diversidade de \textbf{uma} história econômica da agricultura ou por \textbf{uma} história social pelo trabalho escravo criador das riquezas \textbf{que} sustentam o sistema capitalista do mercantilismo ao neoliberalismo, de \textbf{uma} história das sociedades constituídas antes do aparecimento da escrita, da formação de \textbf{uma} civilização americana miscigenada. ( BITTENCOURT, 2018, p.127 ) É \textbf{preciso} \textbf{lembrar} \textbf{que}, apesar de na tradição historiográfica e acadêmica a história factual já estar superada há quase \textbf{um} \textbf{século}, há ainda remanescentes desse factualismo no ensino de História.
É \textbf{preciso} considerar \textbf{que} o professor não é \textbf{um} transmissor de conhecimento e os estudantes, seres passivos \textbf{que} apenas absorvem o saber.
Na BNCC, e mesmo antes dela, com os Parâmetros Curriculares Nacionais ( PCN ), o professor é considerado o mediador do conhecimento e o estudante é \textbf{um} ser ativo, no seu processo de aprendizagem.
Essa forma de aprender e ensinar contribui para a formação do estudante como protagonista.
A aprendizagem de História é \textbf{um} exercício importante de humanização e socialização, pois nos coloca em contato com o outro por meio do conhecimento de outras experiências humanas, em lugares e épocas distintas.
Na BNCC, \textbf{um} dos principais objetivos do componente curricular é estimular a autonomia de pensamento por \textbf{intermédio} do reconhecimento de diferentes \textbf{sujeitos}, histórias, condutas, modos de ser, agir e pensar sobre o mundo.
Tal \textbf{percepção} estimula o pensamento crítico, pois ajuda a compreender \textbf{que} os indivíduos agem de acordo com a época e o lugar nos quais \textbf{vivem}, o \textbf{que} sintetiza \textbf{uma} operação fundamental na construção do conhecimento histórico, qual seja, a \textbf{contextualização}.
Rusen ( 2001 ) corrobora com essa ideia quando afirma \textbf{que} “ a resistência dos \textbf{homens} à perda de si e seu esforço de autoafirmação constituem -se como identidade mediante representações de continuidade, com as quais relacionam as experiências do tempo com as intenções do tempo ” ( p.66 ). \textbf{Um} dos desafios \textbf{que} se coloca no Ensino Fundamental é a necessidade de estudantes e professores assumirem \textbf{uma} “ atitude historiadora ”, dando destaque ao uso das fontes históricas em suas diferentes linguagens¹⁴, realizando progressivas operações cognitivas com as fontes para descrevê -las, analisá -las, compará -las, questioná -las, produzir \textbf{um} discurso sobre o passado e compará -lo com outros discursos já produzidos.
É \textbf{desejável} também ir a campo com os estudantes : observar contextos, entrevistar pessoas, consultar arquivos, bibliotecas, centros de documentação, visitar os lugares de memória, os museus, explorar acervos digitais, coletar e analisar materiais e, por fim, criar seus próprios registros ( como, por exemplo, até mesmo centros de memória na própria escola ). ¹⁴ É também com Bloch e com os Annales \textbf{que} se amplia, já em 1929, a noção de fontes históricas, \textbf{que} passa a ser todo o vestígio deixado pelo \textbf{homem}, desde documentos oficiais a relatos orais.
Atualmente podemos classificar as fontes históricas ou os documentos históricos em : escritas ( cartas, jornais, inventários etc. ) ; orais ( entrevistas, relatos etc. ) ; imateriais ( como festas, saberes \textbf{populares} e tradição oral ) e materiais ( tais como : objetos, construções e indumentária ) ; audiovisuais ( inclusive ficcionais ) ; \textbf{cartográficas} e iconográficas ( pinturas, desenhos, fotografias, entre outros ).
As fontes ou documentos históricos são como pistas de \textbf{um} crime para \textbf{um} detetive ou sintomas de \textbf{uma} doença para \textbf{um} médico como nas analogias apresentadas por Carlo Ginzburg sobre o paradigma indiciário, pois para ele : “ O historiador é, por definição, \textbf{um} investigador para quem as experiências, no sentido rigoroso do termo, estão vedadas ” ( GINZBURG, 1989. p.180 ).
O termo “ atitude historiadora ”, no Currículo Paulista, refere -se ao movimento \textbf{que} professores e estudantes devem realizar para se posicionarem como sujeitos frente ao processo de ensino e aprendizagem, fazendo uso da comparação, \textbf{contextualização} e interpretação das fontes, refletindo \textbf{historicamente} sobre a sociedade na qual \textbf{vivem}, analisando e propondo soluções.
Também, assume -se a atitude historiadora quando se parte do cotidiano do estudante para o passado como \textbf{desdobramento} da “ consciência histórica”¹⁵.
Essa “ consciência ” seria \textbf{inerente} ao ser humano e \textbf{um} resultado das suas interações com o tempo : portanto, o contato de todos com a História se daria antes mesmo do conhecimento sobre os fatos históricos, como decorrência de \textbf{um} processo de existência e sobrevivência humana.
Para Rusen ( 2001 ), > [... ] A consciência histórica é, pois, guiada pela intenção de dominar o tempo \textbf{que} é experimentado pelo \textbf{homem} como ameaça de \textbf{perder} -se na transformação do mundo e dele mesmo.
O pensamento histórico é, por \textbf{conseguinte}, ganho de tempo, e o conhecimento histórico é tempo ganho ( RUSEN, 2001, p.60 ). ¹⁵ Jörn Rüsen, \textbf{pesquisador} alemão, nos trouxe o conceito de “ consciência histórica ” \textbf{que} para ele é o \textbf{fundamento} da ciência histórica.
A aprendizagem de História, \textbf{enquanto} componente da área de Ciências Humanas no Ensino Fundamental, alinha -se às propostas e caminhos do componente de Geografia, o \textbf{que} \textbf{demanda} \textbf{um} trabalho articulado nas escolas, por meio de \textbf{métodos} investigativos em comum e de temáticas semelhantes. (Re)conhecer, identificar, pesquisar, classificar, comparar, diferenciar, interpretar, compreender, analisar, refletir criticamente, criar/ produzir conhecimento a respeito das sociedades humanas em diferentes tempos e espaços, mobilizando várias linguagens ( textuais, iconográficas, \textbf{cartográficas}, materiais, orais, sonoras e audiovisuais ) são propostas dos dois componentes.
O Currículo Paulista propõe \textbf{que} estudantes e professores se coloquem como produtores de conhecimento e \textbf{que} respeitem a diversidade humana.
Desse modo, os estudantes também devem assumir o papel de protagonistas no processo de aprendizagem \textbf{que} tem início nos Anos Iniciais de escolarização e aperfeiçoa -se ao longo da vida, para se tornarem agentes de transformações no meio social.
Todo esse processo contribui para a formação integral¹⁶ do estudante. ¹⁶ Além da BNCC, a educação integral do estudante também \textbf{remete} aos 4 Pilares da Educação, apresentada no relatório da UNESCO para a educação no \textbf{século} XXI, com destaque para \textbf{uma} educação \textbf{que} ultrapassa a aprendizagem sobre técnicas e \textbf{que} define os saberes essenciais : aprender a conhecer, aprender a ser, aprender a conviver e aprender a fazer ( DELORS, 2010 ).
O Organizador Curricular de História está estruturado ano a ano, em unidades temáticas, habilidades e objetos do conhecimento.
O conjunto de habilidades permite o desenvolvimento progressivo das competências específicas de História, da área das Ciências Humanas e das competências gerais da BNCC.
Nos Anos Iniciais do Ensino Fundamental, a escala de observação movimenta -se do particular para o geral.
Assim, no ciclo de alfabetização ( 1º e 2º ano ), propõe -se o estudo do contexto do estudante : o conhecimento de si, do outro, da família, da escola e da comunidade, em continuidade aos saberes desenvolvidos na Educação Infantil, por meio do campo de experiência : “ O eu, o outro, o nós ”.
No 3º ano, amplia -se o objetivo para o estudo da trajetória do município e dos grupos \textbf{que} o formaram¹⁷. ¹⁷ Entendemos \textbf{que} a criança nos anos iniciais do Ensino Fundamental está em processo de alfabetização e \textbf{que} as atividades de registro podem ser realizadas não apenas pela escrita, ou com o auxílio de professores e da família.
No 4º e 5º ano há \textbf{uma} alteração significativa, tendo em vista o \textbf{que} tradicionalmente é aprendido nesta fase, em \textbf{que} a História se desloca do particular e da localidade onde se \textbf{vive} para tempos e espaços mais longínquos.
Tal mudança apresenta -se como possibilidade de melhorar a articulação com os Anos Finais do Ensino Fundamental, diminuindo o descompasso entre essas duas fases da escolarização.
Assim, alguns temas geralmente trabalhados no 6º ano migraram para o 4º e 5º, como o surgimento dos seres humanos e o nomadismo, tendo como ponto de partida o tempo presente marcado por intensos e sucessivos movimentos migratórios.
Outros objetos de conhecimento – como o aparecimento da escrita, da agricultura e de outras tecnologias - também podem garantir esta progressão.

% matched lemmas: cartográfico, conseguinte, contextualização, demandar, desdobramento, desejável, disciplina, enquanto, fundamento, historicamente, homem, inerente, intermédio, lado, lembrar, metodológico, método, percepção, perder, pesquisador, popular, preciso, que, remeter, sujeito, século, um, viver
\end{textsample}
